% GENERATED FILE. Edit canonical lessons, then run npm run generate:books.
% canonical-source-hash: fnv1a64:a5e4c633bcde1c37
% canonical-lessons: IT-C335-impastare, IT-C335-impasto, IT-C335-lievito, IT-C335-lievitare, IT-C335-mattarello

\chapter{Mani in pasta}
\label{ch:it-mani-in-pasta}
\hlchaptermodality{335}

\begin{chapteropening}
\textbf{By the end of this chapter:} \emph{I can describe making dough: kneading, the dough, the yeast, letting it rise and the rolling pin.}

Complete the last lesson of chapter 335: I can describe making dough: kneading, the dough, the yeast, letting it rise and the rolling pin.
\end{chapteropening}

\section[impastare]{impastare — to knead, to mix (a dough)}

\label{lesson:IT-C335-impastare}



\begin{quote}
Before the new one: say the Italian for to deglaze, then the Italian for the soffritto.
\end{quote}

\subsection*{You'll want to know: impastare}
\textbf{impastare} — \textquotedblleft{}to knead, to mix (a dough)\textquotedblright{}.

\begin{quote}
Ho impastato la farina con l'acqua.
\end{quote}

\begin{quote}
\emph{I kneaded the flour with the water.}
\end{quote}

Built on \emph{pasta}, ‘dough’.

\begin{cousinweb}
\emph{impastare} is built on \emph{pasta} \textquotedblleft{}dough, paste\textquotedblright{}, from Late Latin \emph{pasta} \textquotedblleft{}paste\textquotedblright{} (Lewis \& Short), from Greek \emph{pástē}.
\end{cousinweb}

\begin{grammarlens}[title={What you've built}]
One more word for telling, explaining and describing at B1.
\end{grammarlens}

\subsection*{Your turn}
\begin{itemize}
\raggedright
  \item \emph{Say it:} \emph{impastare}
  \item \emph{Say it:} \emph{impastare}, once more
  \item \emph{Say it:} a sentence of your own with \emph{impastare}
  \item \emph{Recall:} say the Italian for to deglaze, then the Italian for the soffritto, then say \emph{impastare} again
\end{itemize}

\begin{tcolorbox}[breakable,colback=teal!4,colframe=teal!35!black,title={Before you move on}]
What does impastare mean? (\textquotedblleft{}To knead, to mix (a dough)\textquotedblright{}.) Say it once more.
\end{tcolorbox}
\section[l'impasto]{l'impasto — the dough, the mixture}

\label{lesson:IT-C335-impasto}



\begin{quote}
Before the new one: say the Italian for the soffritto, then the Italian for to knead.
\end{quote}

\subsection*{You'll want to know: l'impasto}
\textbf{l'impasto} — \textquotedblleft{}the dough, the mixture\textquotedblright{}.

\begin{quote}
L'impasto deve essere morbido.
\end{quote}

\begin{quote}
\emph{The dough must be soft.}
\end{quote}

Any mixture you work with your hands.

\begin{cousinweb}
\emph{impasto} is built on \emph{impastare} \textquotedblleft{}to knead\textquotedblright{}, from \emph{pasta}, Late Latin \emph{pasta} \textquotedblleft{}dough\textquotedblright{} (Lewis \& Short).
\end{cousinweb}

\begin{grammarlens}[title={What you've built}]
One more word for telling, explaining and describing at B1.
\end{grammarlens}

\subsection*{Your turn}
\begin{itemize}
\raggedright
  \item \emph{Say it:} \emph{l'impasto}
  \item \emph{Say it:} \emph{l'impasto}, once more
  \item \emph{Say it:} a sentence of your own with \emph{l'impasto}
  \item \emph{Recall:} say the Italian for the soffritto, then the Italian for to knead, then say \emph{l'impasto} again
\end{itemize}

\begin{tcolorbox}[breakable,colback=teal!4,colframe=teal!35!black,title={Before you move on}]
What does l'impasto mean? (\textquotedblleft{}The dough, the mixture\textquotedblright{}.) Say it once more.
\end{tcolorbox}
\section[il lievito]{il lievito — the yeast, the raising agent}

\label{lesson:IT-C335-lievito}



\begin{quote}
Before the new one: say the Italian for to knead, then the Italian for the dough.
\end{quote}

\subsection*{You'll want to know: il lievito}
\textbf{il lievito} — \textquotedblleft{}the yeast, the raising agent\textquotedblright{}.

\begin{quote}
Aggiungi il lievito e un po' d'acqua.
\end{quote}

\begin{quote}
\emph{Add the yeast and a little water.}
\end{quote}

\emph{Il lievito di birra} — fresh yeast.

\begin{cousinweb}
\emph{lievito} goes back to Latin \emph{levātum}, from \emph{levāre} \textquotedblleft{}to raise\textquotedblright{} (EtymDB).
\end{cousinweb}

\begin{grammarlens}[title={What you've built}]
One more word for telling, explaining and describing at B1.
\end{grammarlens}

\subsection*{Your turn}
\begin{itemize}
\raggedright
  \item \emph{Say it:} \emph{il lievito}
  \item \emph{Say it:} \emph{il lievito}, once more
  \item \emph{Say it:} a sentence of your own with \emph{il lievito}
  \item \emph{Recall:} say the Italian for to knead, then the Italian for the dough, then say \emph{il lievito} again
\end{itemize}

\begin{tcolorbox}[breakable,colback=teal!4,colframe=teal!35!black,title={Before you move on}]
What does il lievito mean? (\textquotedblleft{}The yeast, the raising agent\textquotedblright{}.) Say it once more.
\end{tcolorbox}
\section[lievitare]{lievitare — to rise (of dough)}

\label{lesson:IT-C335-lievitare}



\begin{quote}
Before the new one: say the Italian for the dough, then the Italian for the yeast.
\end{quote}

\subsection*{You'll want to know: lievitare}
\textbf{lievitare} — \textquotedblleft{}to rise (of dough)\textquotedblright{}.

\begin{quote}
Lascia lievitare l'impasto per due ore.
\end{quote}

\begin{quote}
\emph{Let the dough rise for two hours.}
\end{quote}

With \emph{essere}: \emph{è lievitato}. Prices \emph{lievitano} too.

\begin{cousinweb}
\emph{lievitare} is built on \emph{lievito} \textquotedblleft{}yeast\textquotedblright{} (EtymDB), from Latin \emph{levāre} \textquotedblleft{}to raise\textquotedblright{}.
\end{cousinweb}

\begin{grammarlens}[title={What you've built}]
One more word for telling, explaining and describing at B1.
\end{grammarlens}

\subsection*{Your turn}
\begin{itemize}
\raggedright
  \item \emph{Say it:} \emph{lievitare}
  \item \emph{Say it:} \emph{lievitare}, once more
  \item \emph{Say it:} a sentence of your own with \emph{lievitare}
  \item \emph{Recall:} say the Italian for the dough, then the Italian for the yeast, then say \emph{lievitare} again
\end{itemize}

\begin{tcolorbox}[breakable,colback=teal!4,colframe=teal!35!black,title={Before you move on}]
What does lievitare mean? (\textquotedblleft{}To rise (of dough)\textquotedblright{}.) Say it once more.
\end{tcolorbox}
\section[il mattarello]{il mattarello — the rolling pin}

\label{lesson:IT-C335-mattarello}



\begin{quote}
Before the new one: say the Italian for the yeast, then the Italian for to rise.
\end{quote}

\subsection*{You'll want to know: il mattarello}
\textbf{il mattarello} — \textquotedblleft{}the rolling pin\textquotedblright{}.

\begin{quote}
Con il mattarello la pasta diventa sottile.
\end{quote}

\begin{quote}
\emph{With the rolling pin the pastry gets thin.}
\end{quote}

Also written \emph{matterello}.

\begin{cousinweb}
\emph{mattarello} is usually traced to Latin \emph{mateola} \textquotedblleft{}a mallet\textquotedblright{} (Lewis \& Short), but the origin is not settled.
\end{cousinweb}

\begin{grammarlens}[title={What you've built}]
One more word for telling, explaining and describing at B1.
\end{grammarlens}

\subsection*{Your turn}
\begin{itemize}
\raggedright
  \item \emph{Say it:} \emph{il mattarello}
  \item \emph{Say it:} \emph{il mattarello}, once more
  \item \emph{Say it:} a sentence of your own with \emph{il mattarello}
  \item \emph{Recall:} say the Italian for the yeast, then the Italian for to rise, then say \emph{il mattarello} again
\end{itemize}

\begin{tcolorbox}[breakable,colback=teal!4,colframe=teal!35!black,title={Before you move on}]
What does il mattarello mean? (\textquotedblleft{}The rolling pin\textquotedblright{}.) Say it once more.
\end{tcolorbox}
